% =========================================================
% Proof: Carathéodory's Theorem
% Source: volume-iii/analysis/differentiation/notes/notes-caratheodory.tex
% Mode: dual
% =========================================================

\subsection*{Carathéodory's Theorem}
\label{prf:caratheodory}

\begin{remark*}[Return]
\hyperref[thm:caratheodory]{$\leftarrow$ Back to Theorem (Carathéodory's Theorem) in Notes}
\end{remark*}

% ---------------------------------------------------------
% Proof I — Learning Version
% ---------------------------------------------------------
\subsubsection*{Proof I \textemdash\ Learning Version (Bidirectional)}
\label{prf:caratheodory-naive}

\begin{proof}
Let $I \subseteq \mathbb{R}$ be an open interval, let $f : I \to \mathbb{R}$,
and let $c \in I$. The claim is that $f$ is differentiable at $c$ if and only
if there exists a function $\varphi : I \to \mathbb{R}$, continuous at $c$,
satisfying $f(x) - f(c) = \varphi(x)(x - c)$ for all $x \in I$, in which
case $\varphi(c) = f'(c)$.

$(\Rightarrow)$ \textit{Assume $f$ is differentiable at $c$.}

\ProofStep{1}{Define $\varphi$.}
Construct $\varphi : I \to \mathbb{R}$ to represent the difference quotient:
\[
  \varphi(x) :=
  \begin{cases}
    \dfrac{f(x) - f(c)}{x - c} & x \neq c, \\[6pt]
    f'(c)                       & x = c.
  \end{cases}
\]

\ProofStep{2}{Verify the identity.}
For $x \neq c$, multiplying by $(x - c)$ gives $f(x) - f(c) = \varphi(x)(x-c)$
directly. For $x = c$, both sides equal $0$. The identity holds on all of $I$.

\ProofStep{3}{Show continuity at $c$.}
Taking the limit of $\varphi$ as $x \to c$:
\[
  \lim_{x \to c} \varphi(x)
  = \lim_{x \to c} \frac{f(x) - f(c)}{x - c}
  = f'(c)
  = \varphi(c).
\]
Since $\lim_{x \to c} \varphi(x) = \varphi(c)$, $\varphi$ is continuous at
$c$ by \ByDef{Continuity at a Point}.

$(\Leftarrow)$ \textit{Assume there exists $\varphi$ continuous at $c$ with
$f(x) - f(c) = \varphi(x)(x-c)$ for all $x \in I$.}

\ProofStep{4}{Set up the difference quotient.}
For $x \neq c$, divide the identity by $(x - c)$:
\[
  \frac{f(x) - f(c)}{x - c} = \varphi(x).
\]

\ProofStep{5}{Take the limit.}
Since $\varphi$ is continuous at $c$, $\lim_{x \to c} \varphi(x) = \varphi(c)$
by \ByDef{Continuity at a Point}. Therefore:
\[
  \lim_{x \to c} \frac{f(x) - f(c)}{x - c} = \varphi(c).
\]

\ProofStep{6}{Conclude.}
The limit of the difference quotient exists and equals $\varphi(c)$. By
\ByDef{Differentiability at a Point}, $f$ is differentiable at $c$ and
$f'(c) = \varphi(c)$.
\end{proof}

% ---------------------------------------------------------
% Proof II — Professional Standard
% ---------------------------------------------------------
\subsubsection*{Proof II \textemdash\ Professional Standard (Bidirectional)}
\label{prf:caratheodory-direct}

\begin{proof}
Let $I \subseteq \mathbb{R}$ be an open interval, let $f : I \to \mathbb{R}$,
and let $c \in I$. The claim is that $f$ is differentiable at $c$ if and only
if there exists a function $\varphi : I \to \mathbb{R}$, continuous at $c$,
satisfying $f(x) - f(c) = \varphi(x)(x - c)$ for all $x \in I$, in which
case $\varphi(c) = f'(c)$.

$(\Rightarrow)$ \textit{Assume $f$ is differentiable at $c$; construct
$\varphi$ continuous at $c$ satisfying the factorization.}

\ProofStep{1}{Construct $\varphi$.}
Define $\varphi : I \to \mathbb{R}$ by
\[
  \varphi(x) :=
  \begin{cases}
    \dfrac{f(x) - f(c)}{x - c} & x \neq c, \\[6pt]
    f'(c)                       & x = c.
  \end{cases}
\]

\ProofStep{2}{Verify the identity.}
For $x \neq c$, multiplying $\varphi(x) = \frac{f(x)-f(c)}{x-c}$ by $(x-c)$
gives $f(x) - f(c) = \varphi(x)(x-c)$ directly. For $x = c$, both sides
equal $0$. The identity holds on all of $I$.

\ProofStep{3}{Show $\varphi$ is continuous at $c$.}
By \ByDef{Differentiability at a Point},
\[
  \lim_{x \to c} \varphi(x)
  = \lim_{x \to c} \frac{f(x) - f(c)}{x - c}
  = f'(c)
  = \varphi(c).
\]
Since $\lim_{x \to c} \varphi(x) = \varphi(c)$, $\varphi$ is continuous at
$c$ by \ByDef{Continuity at a Point}.

$(\Leftarrow)$ \textit{Assume there exists $\varphi : I \to \mathbb{R}$,
continuous at $c$, with $f(x) - f(c) = \varphi(x)(x-c)$ for all $x \in I$;
deduce that $f$ is differentiable at $c$.}

\ProofStep{4}{Recover the difference quotient.}
For $x \neq c$, dividing the identity $f(x) - f(c) = \varphi(x)(x-c)$ by
$(x - c)$ gives
\[
  \frac{f(x) - f(c)}{x - c} = \varphi(x).
\]

\ProofStep{5}{Pass to the limit.}
By \ByDef{Continuity at a Point} applied to $\varphi$ at $c$,
$\lim_{x \to c} \varphi(x) = \varphi(c)$. Therefore
\[
  \lim_{x \to c} \frac{f(x) - f(c)}{x - c} = \varphi(c).
\]

\ProofStep{6}{Conclude differentiability.}
The limit of the difference quotient exists and equals $\varphi(c)$. By
\ByDef{Differentiability at a Point}, $f$ is differentiable at $c$ and
$f'(c) = \varphi(c)$.
\end{proof}

% ---------------------------------------------------------
% Shared remark blocks
% ---------------------------------------------------------
\begin{remark*}[Proof shape]
The forward direction constructs $\varphi$ explicitly as the difference
quotient extended to $x = c$ by assigning the value $f'(c)$; continuity of
$\varphi$ at $c$ is then immediate from the definition of the derivative.
The converse is a single line: dividing the factorization identity by $(x-c)$
exhibits the difference quotient as $\varphi(x)$, and continuity of $\varphi$
at $c$ supplies the limit.

\FlashProof{1. ($\Rightarrow$) Define $\varphi$ as the difference quotient
  extended by $\varphi(c) = f'(c)$. 2. Verify the identity holds on all of
  $I$. 3. Differentiability gives $\lim_{x\to c}\varphi(x) = \varphi(c)$,
  so $\varphi$ is continuous. 4. ($\Leftarrow$) Divide the identity by
  $(x-c)$ to recover the difference quotient as $\varphi(x)$. 5. Continuity
  of $\varphi$ at $c$ gives the limit $\varphi(c)$. 6. Limit exists, so $f$
  is differentiable with $f'(c) = \varphi(c)$.}
\end{remark*}

\begin{remark*}[Self-assessment of Proof I]
Proof I is correct and the structure is sound. Both directions are properly
set up, citations are present, and the two-case identity verification
(separately for $x \neq c$ and $x = c$) is complete. The following notes
identify the specific differences between the two versions.

\textbf{The two proofs are structurally identical — this is a good sign.}
For Carathéodory's theorem the Learning Version and the Professional Standard
are the same six-step argument. There is no indirect strategy to replace with
a direct one, no unnecessary theorem to eliminate. A student who found this
proof independently has found the professional proof. This will not always
be the case — recognising when it is tells you that your proof instincts on
this kind of construction argument are well-calibrated.

\textbf{Step naming: goal description vs.\ object-and-property.}
Proof I names Step 3 "Show continuity at $c$" — describing the goal as an
action. Proof II names it "Show $\varphi$ is continuous at $c$" — naming the
object ($\varphi$) and its property (continuity at $c$) explicitly. The
second form is preferable: when scanning a proof, the reader should be able
to identify what object is being controlled in each step without reading the
body. Adopt the habit of naming the object first.

\textbf{Step 3: missing citation before the computation.}
Proof I takes the limit $\lim_{x\to c}\varphi(x) = f'(c)$ without naming
the result that licenses this identification. The limit equals $f'(c)$
because that is exactly what \ByDef{Differentiability at a Point} says —
the derivative is defined as that limit. This is a definition invocation and
should be cited before the computation, not left implicit. Proof II adds the
citation explicitly. The rule: every named result used in a computation
appears as a \ByThm or \ByDef citation at the point of use.

\textbf{What Proof I does well.}
The piecewise construction is correct and complete. The $x = c$ case in the
identity verification is correctly checked — this is easy to forget and many
first proofs omit it. The converse direction is concise and correctly
structured. The conclusion in Step 6 cites the right definition. These all
carry forward unchanged into Proof II.
\end{remark*}

\begin{remark*}[Commentary]
\textbf{(a) Why is $\varphi$ defined piecewise?}
The difference quotient $\frac{f(x)-f(c)}{x-c}$ is undefined at $x = c$
since division by zero is not defined. But the theorem requires a function
$\varphi$ defined on all of $I$, including at $c$. The piecewise definition
assigns the only value that can make $\varphi$ continuous: the limit of the
difference quotient as $x \to c$, which is $f'(c)$ by hypothesis.
There is no choice — $\varphi(c) = f'(c)$ is forced by the continuity
requirement.

\textbf{(b) Why check the identity separately at $x = c$?}
For $x \neq c$ the identity holds by algebra. At $x = c$, the right side is
$\varphi(c) \cdot 0 = 0$ and the left side is $f(c) - f(c) = 0$. Both sides
vanish independently of the value of $\varphi(c)$, so the identity holds at
$c$ for free — but it still must be verified explicitly, because the theorem
asserts the identity for all $x \in I$, and $c \in I$.

\textbf{(c) What does the continuity condition actually say?}
Continuity of $\varphi$ at $c$ means $\lim_{x \to c} \varphi(x) = \varphi(c)$.
For $x \neq c$, $\varphi(x)$ is the difference quotient, so
$\lim_{x \to c}\varphi(x)$ is the derivative $f'(c)$ by definition. And
$\varphi(c) = f'(c)$ by construction. The continuity condition and the
differentiability condition are two descriptions of the same limit. This is
why the theorem is an equivalence: neither side adds content that the other
lacks — they are the same statement in different language.

\textbf{(d) Why does the converse divide by $(x-c)$?}
The hypothesis gives $f(x) - f(c) = \varphi(x)(x-c)$ for all $x \in I$.
Dividing by $(x-c)$ for $x \neq c$ produces the difference quotient exactly.
This is pure algebra, not an approximation. The limit is taken separately in
Step 5, where continuity of $\varphi$ supplies the value.

\textbf{(e) Why is continuity of $\varphi$ sufficient to conclude
differentiability?}
After Step 4, the difference quotient equals $\varphi(x)$ for all $x \neq c$.
So $\lim_{x\to c}\frac{f(x)-f(c)}{x-c} = \lim_{x\to c}\varphi(x) = \varphi(c)$
by continuity. The derivative is defined as this limit, so it exists and
equals $\varphi(c)$. The theorem's value is precisely this: it replaces the
problem of showing a limit exists (differentiability) with the problem of
showing a function is continuous at a point (continuity of $\varphi$).
Continuity is often easier to verify algebraically, especially in composed
settings like the Chain Rule.
\end{remark*}

\begin{remark*}[Dependencies]
\begin{itemize}
  \item \hyperref[def:derivative-at-point]{Definition of Differentiability
    at a Point}: Used in Step 3 (Proof II) to identify
    $\lim_{x\to c}\varphi(x)$ as $f'(c)$, and in Step 6 (both proofs) to
    conclude $f$ is differentiable.
  \item \hyperref[def:continuity]{Definition of Continuity at a Point}:
    Used in Step 3 (both proofs) to conclude $\varphi$ is continuous at $c$,
    and in Step 5 (both proofs) to pass the limit through $\varphi$.
\end{itemize>

\FlashUses{def:derivative-at-point, def:continuity}
\end{remark*}